\section{Delimitação do Tema}\label{sec:delimitacao}
Este trabalho delimita-se à análise da metodologia expositiva de ensino aplicada ao contexto do ensino remoto, com foco na avaliação da atenção e do engajamento dos estudantes durante aulas ministradas em ambiente virtual.

A pesquisa será conduzida em ambiente de laboratório, utilizando a ferramenta ESPEON para coleta ativa dos dados comportamentais dos participantes.

O estudo concentra-se na observação de indicadores de tempo e estado das abas, denotadores de foco e interesse discente, obtidos em aulas experimentais ministradas no Centro Federal de Educação Tecnológica Celso Suckow da Fonseca (CEFET/RJ) no mês de novembro de 2025, com um grupo delimitado de alunos, em caráter experimental.

Dessa forma, a investigação não busca avaliar o desempenho acadêmico dos alunos nem comparar metodologias de ensino distintas, mas sim analisar o quão atenta e interessada são as turmas submetidas à metodologia tradicional em sala de aula virtual, com base nos dados empíricos coletados pela ferramenta proposta.